% ch10_c §10.6–10.7 (书页 341–352 / p355–p366)
%--B-TAIL--
%JOIN%
点处出现不连续，紧挨其下方是一个尖锐的峰。这又是向铁磁性转变时观测到的一个特征，尽管理论并不能在所有细节上都重现该峰的确切形状。

铁磁性的巡游电子（itinerant electron）图像对铁磁金属绝大多数观测性质的解释，与局域自旋模型同样成功。它也与普遍的能带理论（§3.10）、输运理论以及费米面研究相一致；例如，在其中可以观察到两套分立的能带系统，分别容纳“向上”自旋与“向下”自旋的电子。另一方面，“交换”能 $U$ 的确切本性一点也不清楚，而且有独立的证据表明磁矩存在某种程度的表观局域化（§10.6、§10.11）。当前最好的看法是：对具有 d 带的过渡金属倾向于能带模型，而稀土金属中部分填充的 f 态里的自旋，则似乎总是表现得如同按 Heisenberg 方式局域化。

如 §8.3 所述，有些物质是铁电体（ferro-electric）：它们在没有电场的情况下也能携带永久电矩。这一现象在宏观上与铁磁性相似，但其来源是一个复杂得多、也极不寻常的微观效应。它出自某些晶体结构的一种能力：使自己发生轻微畸变，从而让晶格的每一个晶胞都带有一个偶极矩。

\section{磁性杂质}

过渡元素或稀土元素的原子溶入普通金属后，往往保留局域磁矩。这是一系列有趣的实验现象和理论概念的源泉。用相移语言（§5.5）所作的 Friedel 解释，与通过 Anderson 哈密顿量（Anderson Hamiltonian）的传统描述互为补充；借助后者，可以用少数几个唯象参量对这类现象的许多特征作出半定量估计。

首先我们假设：自旋相反的电子之间的“交换”或“关联”排斥，只有当二者处在同一个能量为 $\mathcal{E}_d$ 的原子 $d$ 能级上时才起作用。利用式 (10.38) 中已经用过的简化 Hartree--Fock 近似，我们写出

\begin{equation*}
\mathcal{H}_{dd} = U n_{d+} n_{d-} \tag{10.44}
\end{equation*}
其中 $n_{d\sigma}$ 表示自旋为 $\sigma$ 的 $d$ 电子的平均数目。倘若过渡元素溶入普通金属后这些原子能级仍保持完全锐利，它们无疑会被极化——比方说 $n_{d+}$ 取尽可能大的值，而 $n_{d-}$ 取尽可能小的值——这正是 §10.5 所讨论的那类情形的一个夸张版本。在这种组态下，各能级的有效能量被劈开，即（图 178）

\begin{equation*}
\mathcal{E}_{d\sigma} = \mathcal{E}_d + U n_{d\sigma}. \tag{10.45}
\end{equation*}
%FIG{178}: {原子 $d$ 能级（a）被 $d$–$d$ 相互作用极化并劈裂（b），随后因与 $s$ 电子杂化而展宽为共振峰（c）。}

现在我们在这些局域 $d$ 态与溶剂金属传导电子（“s 电子”）的 Bloch 态之间引入一个强度为 $V$ 的相互作用。这个相互作用不一定对自旋敏感；原则上，它与过渡金属能带结构的模型哈密顿量表述（§3.10）中的 $s$–$d$ 杂化矩阵元基本相同，而那个矩阵元本可以由诸如式 (3.27) 那样的 L.C.A.O.“重叠积分”导出。因此我们假设这个相互作用只联系同种自旋的 $s$ 态和 $d$ 态，并且忽略 $s$ 电子自身之间可能存在的微弱交换相互作用。

然而，若把这个相互作用看作作用在每个 $d$ 能级上的微扰，它将引起与 $s$ 带之间的往返跃迁，其速率为

\begin{equation*}
W_\sigma/\hbar = (\pi/\hbar)\, V^2 \mathcal{N}_s(\mathcal{E}_{d\sigma}) \tag{10.46}
\end{equation*}

在这个“黄金定则”（Golden Rule）公式里，溶剂金属的性质仅通过其 $s$ 带中处于所考虑 $d$ 态能量 (10.45) 处的态密度体现出来。但这段被缩短的寿命可以解释为 $d$ 能级在能量上展宽了 $W_\sigma$。在 $\mathcal{E}_{d\sigma}$ 处的锐利能级不复存在，我们得到的是典型 Lorentz 线型的谱密度函数

\begin{equation*}
\rho_{d\sigma}(\mathcal{E}) = \frac{2l+1}{\pi}\, \frac{W_\sigma}{(\mathcal{E} - \mathcal{E}_{d\sigma})^2 + W_\sigma^2}. \tag{10.47}
\end{equation*}

原则上，这条谱线的中心也会稍有移动，但这不是一个重要的物理效应。

这个公式其实正是我们把 Friedel 求和 (5.40) 作为能量的函数、在诸如 (3.81) 那样的共振附近求导时应得的结果。因此，把宽度 $W_\sigma$ 解释为 (10.46) 中 $s$–$d$ 矩阵元 $V$ 的做法，并不是物理本质的要素，尽管在这些相互作用不太强时这些方程是有效的。

但是现在，一旦杂质上的 $d$ 能级被展宽，我们就不能再肯定它们仍会保持磁性极化。此时的情形其实与 (10.39)–(10.42) 完全一样，只需用 (10.47) 代替 $\frac{1}{2}\mathcal{N}(\mathcal{E})$。于是我们得到杂质离子上存在永久矩的一个条件：如同 (10.43)，若

\begin{equation*}
U \rho_{d+}(\mathcal{E}_F) > 1 \tag{10.48}
\end{equation*}
则极化态在能量上就是稳定的，不会有电子从 $n_{d+}$ 分布转移到 $n_{d-}$ 分布（参见 §5.2）。

过渡元素和稀土元素稀薄合金的多种可观测性质，都可以由这个理论得出，其中的参量 $U$ 和 $V$ 取合理而自洽的数值。例如，从 V 到 Co 的铁族元素溶解在 Cu、Ag、Au 中时是磁性的，但作为杂质出现在 Al 中时却不显示磁矩。Al 中传导电子浓度较高，使 (10.46) 中的 $\mathcal{N}_s$ 过大，从而使 $d$ 能级展宽，压低 (10.47) 中 $\rho_{d\sigma}(\mathcal{E})$ 的峰值，以致磁化条件 (10.48) 不再能得到满足。

基体金属的传导电子对杂质的存在并非无动于衷，只是这一效应依赖于自旋。设磁性离子完全极化，总自旋为 $S$。若所有“向上”自旋的 $d$ 态都已填满，同种自旋的 $s$ 电子便无法进入。因此这时只有“向下”自旋的 $s$ 电子受益。按普通微扰论估计，经由矩阵元 $V$ 虚跃迁到比费米能级高出能量 $U$ 的 $d$ 能级上，会使这些电子中每一个的能量改变

\begin{equation*}
J/N \approx -V^2/U. \tag{10.49}
\end{equation*}
于是系统的行为就好像存在如下形式的依赖于自旋的相互作用能

\begin{equation*}
\mathcal{H}_{sd} = -(J/N)\, \mathbf{s} \cdot \boldsymbol{\mathsf{S}}_i\, \delta(\mathbf{r} - \mathbf{R}_i) \tag{10.50}
\end{equation*}
杂质在 $\mathbf{R}_i$ 处的局域自旋矩 $\boldsymbol{\mathsf{S}}_i$ 与传导电子在 $\mathbf{r}$ 处的自旋 $\mathbf{s}$ 联系起来。注意，无论 $V$ 的符号如何，这一相互作用的符号都倾向于使两个自旋反平行（见 §10.7）。

这个 $s$–$d$ 相互作用 (10.50) 是作用在自由电子气上的一个高度局域化的微扰。为了领会它的效应，让我们计算系统的磁“响应”作为波数的函数，做法与我们在 §5.1 中计算静电响应函数基本相同。眼前的问题其实要简单得多，因为我们忽略了电子气内部的自旋–自旋相互作用，可以把整个效应都归因于“外”微扰。“向上”自旋和“向下”自旋的 $s$ 电子各带磁矩 $g\beta\mathbf{s}$，可以分开处理；按照导出 (5.16) 的论证，我们得到一个广义磁化率（参见 §8.1）：

\begin{equation*}
\chi_0(q) = 2(g\beta s)^2 \sum_{\mathbf{k}} \frac{f^0(\mathbf{k}) - f^0(\mathbf{k}+\mathbf{q})}{\mathcal{E}(\mathbf{k}+\mathbf{q}) - \mathcal{E}(\mathbf{k})}, \tag{10.51}
\end{equation*}
其中求和遍及两种自旋的电子态。

在自由电子模型中，这个和可以像 (5.35) 那样求出，结果是一个简单的解析函数，在 $q = 2k_F$ 处有二阶导数奇性。再作 Fourier 变换（参见式 (5.17)），便给出局域矩 (10.50) 邻域内磁化强度的空间分布。该奇性产生一个波数为 $2k_F$ 的振荡项，在远距离处按 $1/r^3$ 衰减（图 179）。这就是 R.K.K.Y.（Ruderman--Kittel--Kasuya--Yosida）效应；距杂质远近不同的各个离子，其 Knight 位移（§10.2）会有微小变化，可以用它来探测这一效应。

用 §5.5 的语言，我们也可以把磁性极化的杂质描述成一个区域：“向上”自旋电子因 d 相移中出现共振 (3.81) 而被吸引进去。而“向下”自旋电子的相应共振出现在高得多的能量处，因此两种分布在费米能级处的相移相差很大。由 (5.41)，这就给出磁极化的振荡，与静电杂质周围的电荷 Friedel 振荡完全类似。这种描述与 Anderson 模型等价，但不能像 (10.49) 那样直接给出用 $V$ 和 $U$ 表示的 $J$ 的数值。

或许更重要的是，这一现象为两个相邻磁性离子上磁矩之间的耦合提供了一种机制。“向上”自旋杂质紧邻处“向下”自旋 $s$ 电子的过剩，会在该邻域内的任何其他杂质上诱导出相应的“向上”自旋。

%FIG{179}: {$s$ 电子围绕杂质 A 的自旋极化 Friedel 振荡，导致与近邻杂质 B 的铁磁性相互作用，但在更远的距离（如 C 处）则可能有利于反铁磁性。}

这样，我们就有了铁磁性的 Heisenberg 交换模型（§10.4）所需的全部机制：过渡金属离子 $d$ 壳层中局域化的磁矩，由于其极化方向倾向于与浸泡着它们的传导电子气反平行，而被促成为彼此平行。

这类相互作用在稀薄合金中引起复杂的磁学和热学效应，其中磁性离子是无规分布在溶剂晶格上的。在极端情形下，把非磁性溶剂去掉，留下纯过渡金属，这就是诸如 Fe 等金属中铁磁性的 Vonsovski--Zener 模型（Vonsovski--Zener model）。注意，这是一个比 Heisenberg--Dirac 哈密顿量 (10.29) 所提示的“动力学”得多的模型。$d$ 壳层磁矩并不是由整数个完全局域的“d 电子”组成，而是代表 $s$ 带与 $d$ 带已充分杂化的 Bloch 态中的平均贡献。因此，这个模型与 §10.5 中讨论的较简单的能带铁磁性并不矛盾，只是参量 $U$ 需要重新解释。

依赖于自旋的 $s$–$d$ 相互作用 (10.50) 还导致一个与磁性杂质对电阻率的贡献相联系的奇特效应。设我们像 (6.90) 中那样，试图计算这样一个物体把一个处于态 $|\mathbf{k},+\rangle$ 的“向上”自旋 $s$ 电子散射到同自旋态 $|\mathbf{k}',+\rangle$ 的微分截面。

在第一级 Born 近似下，散射振幅恰好就是矩阵元

\begin{align*}
t^{(1)} &= \langle \mathbf{k}, + | \mathcal{H}_{sd} | \mathbf{k}', + \rangle \notag \\
&= -(J/N)\, \mathsf{S}_i^{(z)}. \tag{10.52}
\end{align*}
这一项平方后不为零，而且是价电荷散射等（§7.4）之外的附加项，本来通常就足以满足我们对此问题的好奇心了。

%FIG{180}: {二级 $s$–$d$ 散射中的直接过程（a）与交换过程（b）。注意自旋翻转跃迁中 $s$ 自旋与 $d$ 自旋的相对取向。}

但是，本着一种迂腐的精神，让我们考察微扰级数中更高阶的项。初等量子力学教科书教导我们，第二级 Born 近似取如下形式

\begin{equation*}
t^{(2)} = \sum_{\mathbf{k}'',\sigma} \frac{\langle \mathbf{k}, + | \mathcal{H}_{sd} | \mathbf{k}'', \sigma \rangle \langle \mathbf{k}'', \sigma | \mathcal{H}_{sd} | \mathbf{k}', + \rangle}{\mathcal{E}(\mathbf{k}) - \mathcal{E}(\mathbf{k}'')}, \tag{10.53}
\end{equation*}
其中“中间”态 $|\mathbf{k}'',\sigma\rangle$ 通过相互作用哈密顿量 $\mathcal{H}_{sd}$ 与初态和末态两者相联系。然而，这个表达式并不完全正确。在完整的时间依赖微扰论中，我们必须考虑两类“图”（图 180）。在直接
过程（a）中，入射电子先进入态 $|\mathbf{k}'',\sigma\rangle$（因此该态必须是空的），然后才被发射到末态；而在交换（exchange）过程（b）中，我们假设一个已处于 $|\mathbf{k}'',\sigma\rangle$ 的电子，在“初始”电子被接纳进中间态之前，就先以 $|\mathbf{k}',+\rangle$ 离去。由这些规则，我们应当把 (10.53) 改写成

\begin{align*}
t^{(2)} = \sum_{\mathbf{k}'',\sigma} \frac{1}{\mathcal{E}(\mathbf{k}) - \mathcal{E}(\mathbf{k}'')} \big\{ &(1 - f_{\mathbf{k}''}) \langle \mathbf{k}, + | \mathcal{H}_{sd} | \mathbf{k}'', \sigma \rangle \langle \mathbf{k}'', \sigma | \mathcal{H}_{sd} | \mathbf{k}', + \rangle \\
&+ f_{\mathbf{k}''} \langle \mathbf{k}'', \sigma | \mathcal{H}_{sd} | \mathbf{k}', + \rangle \langle \mathbf{k}, + | \mathcal{H}_{sd} | \mathbf{k}'', \sigma \rangle \big\}, \tag{10.54}
\end{align*}
其中 $f_{\mathbf{k}''}$ 是中间态的占据数。

对于普通的势散射，(10.54) 中两个矩阵元的乘积与其次序无关，于是 $f_{\mathbf{k}''}$ 相消，我们便回到 (10.53)，它对散射截面的贡献只是对一级项 (10.52) 的一个小的修正。但对自旋散射我们必须记住，$\mathcal{H}_{sd}$ 中含有诸如 $\mathsf{S}_i^{(x)} s_x$ 和 $\mathsf{S}_i^{(y)} s_y$ 这样的项，它们以消耗杂质的极化为代价去“翻转”传导电子的自旋（见 (10.104)）。算符 $\mathsf{S}_i^{(x)}$ 与 $\mathsf{S}_i^{(y)}$ 不对易，因此 (10.54) 中各项的先后次序变得重要。从杂质自旋矩的角度看，两个跃迁中哪一个先发生确实至关重要。例如显而易见，倘若杂质已经具有最大的“向上”自旋，直接过程就是被禁戒的。

结果是，在 (10.54) 的代数约化中最终留下一个如下形式的项

\begin{equation*}
t_{\rm K} \approx 2(J/N)^2 \mathsf{S}_i^{(z)} \sum_{\mathbf{k}''} \frac{f_{\mathbf{k}''}}{\mathcal{E}(\mathbf{k}) - \mathcal{E}(\mathbf{k}'')}, \tag{10.55}
\end{equation*}
其中因子 $f_{\mathbf{k}''}$ 并未被来自其他图的贡献所抵消。由于 $f_{\mathbf{k}''}$ 在费米能量处陡然变化，当 $\mathcal{E}(\mathbf{k})$ 趋近 $\mathcal{E}_F$ 时，这个和会变得很大。例如，假设传导带中的态密度一直低到某个能量 $(\mathcal{E}_F - D)$ 都具有 $\mathcal{N}(\mathcal{E}_F)$ 的量级，我们可以把 (10.55) 近似地算成

\begin{equation*}
t_{\rm K} \approx 2(J/N)^2 \mathsf{S}_i^{(z)} \mathcal{N}(\mathcal{E}_F) \ln \frac{D}{|\mathcal{E}_F - \mathcal{E}(\mathbf{k})|}. \tag{10.56}
\end{equation*}
普通带电杂质在金属中产生的\emph{剩余电阻}（residual resistance）与温度无关（§§7.2--7.4）。但是，当我们把磁性杂质引起散射的总概率在费米面上厚度为 $kT$ 的“热层”内积分时，
(10.56) 中的奇性并未因此被抹平。电阻率应当随温度作对数式变化，大致为

\begin{equation*}
\rho(T) \sim \rho_0 \{1 - 2(J/N)\,\mathcal{N}(\mathcal{E}_F) \ln D/kT\}. \tag{10.57}
\end{equation*}

由于 $J$ 本质上是负的，这个公式预言剩余电阻率随温度下降而迅速增大。再把它与基体金属低温下按 $T^5$ 变化的“理想”电阻（§7.5）结合起来，就解释了长期在含过渡元素的稀薄合金中观测到的神秘的电阻极小（resistance minimum）。而 (10.56) 的强能量依赖进入诸如 (7.104) 那样的公式，则解释了在这类材料中同样观测到的巨大温差电功率（giant thermo-electric power）。

但完整的 Kondo 效应（Kondo effect）理论还必须考虑近邻杂质的相互极化、$t$ 的微扰展开的高阶修正、具有 $s$–$d$ 相互作用的系统之基态的性质，以及许多其他微妙之处。这个领域如今已成为固体物理中各种高深量子力学方法的主要演练场之一。

%FIG{181}: {（a）铁磁有序。（b）反铁磁有序。}

\section{反铁磁性}

Heisenberg 模型中铁磁性的一般条件是 (10.29) 中的交换积分 $J_{ll'}$ 为正。此时内场与自旋的排列方向平行，整个系统倾向于取图 181(a) 的组态。

但假如——正如对 $J_{ll'}$ 的计算常常表明的那样——该积分的符号为负，那么任一格点上的内场 (10.32) 就会与近邻自旋磁化的平均方向相反，图 181(b) 的排列成为有利。我们把这样的系统称为反铁磁性（antiferromagnetic）的。

为了看出会发生什么，让我们分别考虑自旋取向相反的两套子晶格。设
$\boldsymbol{\mu}_+$ 是 $\oplus$ 子晶格上原子的平均磁化。交换相互作用将在 $\ominus$ 子晶格上产生一个内场（不妨写成）

\begin{align*}
\mathbf{H}^- &= \frac{1}{\beta^2}(\Sigma J)\boldsymbol{\mu}_+ \notag \\
&= \lambda N \boldsymbol{\mu}_+, \tag{10.58}
\end{align*}
这个求和假定只对最近邻进行，因为 $J_{ll'}$ 被认为作用程极短。要点在于，$\lambda$ 现在将是\emph{负}的。

同样由于这一相互作用，每个 $\oplus$ 格点上也存在内场

\begin{equation*}
\mathbf{H}^+ = \lambda N \boldsymbol{\mu}_-. \tag{10.59}
\end{equation*}

现在假设我们的磁体是经典的，即 (10.17) 对磁场 $\mathbf{H}$ 中磁矩 $\boldsymbol{\mu}$ 的\emph{平均}值成立。把 (10.17) 写成

\begin{equation*}
\langle \boldsymbol{\mu} \rangle = \mathcal{L}\{(\boldsymbol{\mu}\cdot\mathbf{H})/kT\}, \tag{10.60}
\end{equation*}
其中用符号 $\mathcal{L}$ 表示由 (10.17) 右端定义的 Langevin 函数（Langevin function）。于是我们有两个方程要求解：

\begin{align*}
\boldsymbol{\mu}_- &= \mathcal{L}\{\boldsymbol{\mu}\cdot(\mathbf{H}+\lambda N\boldsymbol{\mu}_+)/kT\}, \notag \\
\boldsymbol{\mu}_+ &= \mathcal{L}\{\boldsymbol{\mu}\cdot(\mathbf{H}+\lambda N\boldsymbol{\mu}_-)/kT\}. \tag{10.61}
\end{align*}

在高温下，当近似 (10.18) 有效时，这两个方程可以相加而给出 (10.23)。结果与先前的 (10.24) 几乎完全一样：

\begin{equation*}
\chi = \frac{N\langle\mu^2\rangle}{3k(T+\theta)}; \tag{10.62}
\end{equation*}
$\lambda$ 为负时，方便的做法是按 $-\lambda N\langle\mu^2\rangle/3kT$ 来定义 $\theta$，这与 (10.25) 相一致。磁矩倾向于彼此反平行地排列，而不是被磁场拉向同一方向，顺磁磁化率因此减小。

但在低温下，即使不存在磁场也有一个解。若我们写出

\begin{equation*}
\boldsymbol{\mu}_- = -\boldsymbol{\mu}_+ \tag{10.63}
\end{equation*}
（在图 181(b) 的有序反铁磁排列中，直觉上正是如此），则 $\lambda$ 的负号被抵消，我们得到方程

\begin{equation*}
\boldsymbol{\mu}_+ = \mathcal{L}\{(\boldsymbol{\mu}\cdot|\lambda|N\boldsymbol{\mu}_+)/kT\} \tag{10.64}
\end{equation*}
待解。这个方程与 §10.3 中导致铁磁性的方程完全相同——例如 (10.26)，其中含有 Langevin 函数的“量子”等价物。

这样，从单个自旋的角度看，铁磁性与反铁磁性似乎很相似：系统在高温下是顺磁性的，但在某一温度以下变成强极化；在反铁磁情形中，这个温度称为 Néel 温度（Néel temperature），

\begin{equation*}
T_{\rm N} = \theta. \tag{10.65}
\end{equation*}

然而在宏观上，两者的行为大不相同。反铁磁排列不具有任何净极化，因此系统不显示任何永久磁化。所能观测到的只是磁化率的变化；要计算它，需在一个小的磁场 $H$ 之下求解 (10.61)。

%FIG{182}: {反铁磁体的磁化率：（a）$\chi_\parallel$ 趋于零；（b）$\chi_\perp$ 保持恒定；（c）$\chi = \frac{1}{3}\chi_\parallel + \frac{2}{3}\chi_\perp$。}

这个计算原则上很简单，只是实际做起来略嫌繁琐。其结果有趣之处在于，它依赖于 $\mathbf{H}$ 相对于两套子晶格磁化方向的取向。当 $\mathbf{H}$ 与 $\boldsymbol{\mu}_+$ 平行时，磁化率 $\chi_\parallel$ 在 $T=0$ 处趋于零。这是因为内场强到足以阻止 $\ominus$ 子晶格中的任何自旋翻转。另一方面，$\chi_\perp$ 在 $T=0$ 的值与它在 $T_{\rm N}$ 处应有的值相同：与外场垂直的子晶格极化，对其磁化被稍微转向外场方向只有微弱的抗拒。

实际上，一个宏观样品将包含大量磁畴，各磁畴的极化轴彼此不同，于是对平均磁化率可以写出

\begin{equation*}
\chi = \tfrac{1}{3}\chi_\parallel + \tfrac{2}{3}\chi_\perp. \tag{10.66}
\end{equation*}
$\chi$ 在 $T_{\rm N}$ 处斜率的不连续性很容易观测到，并且与子晶格极化时出现的比热反常相联系。

值得指出的是，由于 $\chi_\perp > \chi_\parallel$，磁场本应能够把所有磁畴都拉成垂直组态。但这种情况并不发生——大概是由于自旋–自旋相互作用的各向异性，它倾向于偏爱某些特殊的晶体方向。

我们在这里考虑的模型，只是展现着令人眼花缭乱的种种现象的一大类系统中的一个简单例子。这个论证可以从不同方向加以细化。例如，晶体结构可能是复杂的，其中磁性离子并不形成简单立方网络，系统究竟会如何划分成子晶格也可能并非一目了然。又如，相互作用的作用程可能比最近邻间更长，或者对某些自旋对是反铁磁性的，而对另一些则是铁磁性的。这样做的一个后果（在我们的框架内很容易处理：只需在 (10.58) 中加上诸如 $\lambda' N\boldsymbol{\mu}_-$ 的项，在 (10.59) 中加上 $\lambda' N\boldsymbol{\mu}_+$ 的项）是改变 Néel 温度，使之不再与 (10.62) 中的参量 $\theta$ 相同。

%FIG{183}: {亚铁磁性。}

%FIG{184}: {超交换。}

另一种情形是晶体中存在两种不同类型的磁性离子，而且二者不具有相同的磁矩。于是我们可能遇到图 183 所描绘的情况：A、B 两类原子的自旋取向相反，但由于一套子晶格上的矩比另一套上的大，系统仍将显示永久矩。这称为亚铁磁性（ferrimagnetism），是铁氧体（ferrite）的特征；在宏观上，它看起来与铁磁性非常相像。

在许多反铁磁物质中，磁性离子彼此相距相当远，以致很难相信它们能由诸如 (10.37) 那样的交换积分联系起来。这并不意味着晶格非常稀疏，而是意味着在磁性离子之间的空隙里填有许多非磁性离子。例如，它们可以是氧离子，如 MnO 中那样。这时我们就遇到了称为超交换（super-exchange）的机制：Mn 离子中的 $d$ 电子与氧中的电子相耦合，而后者再通过交换作用与第二个 Mn 离子中的 $d$ 电子相互作用。

示意性地，情形可以如图 184 所示。若原子 A 具有 $+$ 自旋，就有空位让一个 $-$ 自旋的电子从氧的 $p$ 态进入。若原子 B 具有 $-$ 自旋，它就能接纳这对电子中来自氧的另一个。这样，氧上 $p$ 电子自旋之间的相互作用就转化为磁性离子 $d$ 壳层自旋之间的相互作用。这一效应的实际计算颇为微妙，但它至少在原理上解释了观测到的现象。

在稀磁合金中，§10.6 的间接交换（indirect exchange）机制也可以是反铁磁性的。如图 179 所示，当第二个杂质恰好处在与第一个磁性离子的 R.K.K.Y. 振荡中“向上”自旋峰相对应的距离上时，就会出现这种情形。但要理解诸如 Cr 这类纯金属中的反铁磁性，我们应当回到 §10.5 的巡游电子图像。

定量地看，不难理解反铁磁图样如何能被传导电子之间的交换作用所稳定。设图 181(b) 那种类型的磁有序已经建立。一个“向上”自旋的电子穿越晶体时，通过交换相互作用 $U$ 将感受到一个依赖于局域极化的有效场。这个场具有磁超晶格的周期性，因而会在通常布里渊区内部的一个“子区”边界上引起能量不连续（参见 §2.6）。能带结构中的这一特征会反映为“向上”自旋电子密度的一个超晶格。对于自旋相反的电子，微扰的符号相反，于是它们倾向于占据另一套超晶格。如果这些分布碰巧与我们最初设想的磁化图样相一致，我们就找到了电子的一种自洽的反铁磁有序，而无须强迫它们局域化。

为了给这类论证赋予数学实质，我们假设
